%% OrdinalFed talk -- technical speech + study guide (learn the terms behind the simple speech)
\documentclass[11pt,a4paper]{article}
\usepackage[margin=2.1cm]{geometry}
\usepackage[T1]{fontenc}
\usepackage[scaled]{helvet}
\renewcommand{\familydefault}{\sfdefault}
\usepackage{amsmath,amssymb}
\usepackage{xcolor}
\usepackage{enumitem}
\usepackage{longtable,booktabs,array}
\usepackage{needspace}
\usepackage[hidelinks,pdftitle={OrdinalFed - Technical Speech and Study Guide}]{hyperref}
\definecolor{navy}{HTML}{14213D}
\definecolor{accent}{HTML}{2F6FB3}
\definecolor{muted}{HTML}{5B6575}
\definecolor{termbg}{HTML}{EEF3F9}
\setlength{\parindent}{0pt}
\setlength{\parskip}{5pt}
\linespread{1.2}
\pagestyle{plain}

\newcommand{\slide}[3]{%
  \par\needspace{7\baselineskip}\bigskip\noindent{\color{accent}\large\bfseries Slide #1 \textperiodcentered\ #2}%
  \hfill{\color{muted}\small #3}\par\smallskip}
\newcommand{\act}[1]{{\color{muted}\itshape (#1)}}
% "Terms on this slide" list
\newenvironment{terms}{%
  \par\smallskip\noindent{\color{accent}\small\bfseries TERMS ON THIS SLIDE}\par
  \begin{itemize}[leftmargin=1.2em,itemsep=2pt,topsep=2pt,label={\color{accent}\textbullet}]\small}%
  {\end{itemize}}
\newcommand{\term}[1]{\textbf{#1:}}

\begin{document}

{\color{navy}\LARGE\bfseries OrdinalFed: Technical Speech \& Study Guide}\par\smallskip
{\color{muted}IEEE WIECON-ECE 2026 \textperiodcentered\ Paper ID 518 \textperiodcentered\ Md Tamim Hossain}

\medskip
\fcolorbox{accent}{termbg}{\parbox{\dimexpr\linewidth-2\fboxsep-2\fboxrule}{%
\textbf{How to use this guide.} This is the \emph{technical} version of the talk, with the proper
terms. After each slide, the blue \textbf{Terms on this slide} list explains every technical word in
easy language. Use it to learn, so you can answer questions. On stage, read the \emph{simple}
speech (\texttt{OrdinalFed\_speech.pdf}). At the end: a quick A--Z glossary and harder questions
with answers.}}

%% ------------------------------------------------------------------
\slide{1}{Title}{0:20}
Good morning, everyone. I am Md Tamim Hossain from the Department of Computer Science and
Engineering, MIST. Today I will present our paper, ``OrdinalFed: Federated Learning with
Ordinal-Aware Aggregation for Diabetic Retinopathy Grading.'' This is joint work with Najifa
Raksanda, M.\,M. Tamim Sharif and Aria Nawshin Ashka.
\begin{terms}
\item \term{Federated learning (FL)} many hospitals train \emph{one} AI model together, but each
hospital keeps its patient images at home. Only the model travels, never the images.
\item \term{Ordinal} things that have a natural order, like small $<$ medium $<$ large. DR grades
0--4 are ordinal: Grade 3 is between Grade 2 and Grade 4.
\item \term{Aggregation} the step where the server combines the hospitals' models into one model.
``Ordinal-aware aggregation'' means the combining step pays attention to the order of the grades.
\item \term{Grading} giving each eye image a severity level (a grade), not just ``sick or healthy''.
\end{terms}

%% ------------------------------------------------------------------
\slide{2}{Diabetic retinopathy}{0:45}
Diabetic retinopathy, or DR, is damage to the small blood vessels of the retina caused by
diabetes. It is one of the main causes of preventable blindness. In Bangladesh alone, more than
13 million people have diabetes, and outside the big cities there are very few eye specialists.
Doctors grade DR on five levels, from Grade 0, no disease, to Grade 4, the most severe. Finding
Grade 2 and above early can prevent vision loss. So an automatic grading model would be very
useful.
\begin{terms}
\item \term{Retina} the thin layer at the back of the eye that senses light, like the sensor in a
camera.
\item \term{Fundus photograph} a photo of the back of the eye (the retina) taken with a special
camera. Our datasets are made of these photos.
\item \term{The five grades} 0 = no DR; 1 = mild; 2 = moderate; 3 = severe; 4 = proliferative.
Grades 1--3 are called \emph{NPDR} (non-proliferative DR). Grade 4 is \emph{PDR}: new, weak blood
vessels start to grow, and they can bleed and cause blindness.
\item \term{Signs in the photo} \emph{microaneurysms} (tiny bulges in vessels), \emph{haemorrhages}
(small bleeds), \emph{hard exudates} (yellow fat/protein spots), and new vessels in Grade 4.
\item \term{Referable DR} Grade 2 or higher. These patients should be sent (``referred'') to an eye
specialist.
\end{terms}

%% ------------------------------------------------------------------
\slide{3}{Why federated learning?}{0:45}
Deep learning models grade DR well when all images are collected in one place. But privacy rules
make it hard for hospitals to share patient images. Federated learning solves this. Each hospital
trains the model on its own images. Then it sends only the model weights to a server, never the
images. The server combines them into one global model and sends it back. This repeats for several
rounds. The standard way to combine the models is called FedAvg. It weights each hospital by how
many images it has.
\begin{terms}
\item \term{Deep learning} AI that uses large neural networks, which learn patterns from many
examples (here, eye photos with known grades).
\item \term{Model weights (parameters)} the millions of numbers inside a neural network. Training
means adjusting these numbers. In FL, hospitals share these numbers, not images.
\item \term{Client} one participant in FL, here one hospital. \term{Server} the central computer
that combines the clients' models.
\item \term{Global model} the combined model made by the server, shared with every hospital.
\item \term{Communication round} one full cycle: server sends the global model $\rightarrow$
hospitals train $\rightarrow$ hospitals send models back $\rightarrow$ server combines. We use 10
rounds.
\item \term{FedAvg (Federated Averaging)} the standard FL method (McMahan et al., 2017). The server
takes a \emph{weighted average} of the hospitals' models; a hospital's weight is its share of all
images. Example: a hospital with 60\% of the images gets 60\% of the say.
\item \term{Privacy rules (HIPAA, GDPR)} laws in the USA and Europe that protect patient data.
\end{terms}

%% ------------------------------------------------------------------
\slide{4}{Two problems}{0:50}
There are two problems. First, hospitals see very different patients. A screening clinic sees
mostly healthy eyes, while a referral hospital sees many severe cases. When the data are this
uneven, FedAvg can collapse. In our experiments under extreme skew, its test QWK dropped from
0.647 in its best round to 0.378 in its worst. QWK is the standard score for ordered grades: 1
means perfect, and far-off mistakes cost more.

Second, the usual loss, cross-entropy, ignores the order of the grades. For cross-entropy, calling
a Grade 4 image Grade 0 is no worse than calling it Grade 3. But clinically, the first mistake is
much worse.
\begin{terms}
\item \term{Non-IID data} ``IID'' means every hospital's data looks like the overall data.
\emph{Non-IID} means each hospital's data is different. Here the difference is in the grades, which
is called \emph{label skew}: one hospital has mostly Grade 0, another mostly Grade 2.
\item \term{Extreme skew} our hardest test, where each hospital sees mostly one grade (made with a
Dirichlet split, $\alpha = 0.1$; see Slide 10).
\item \term{Collapse} the global model suddenly gets much worse in later rounds. Cause:
\emph{client drift}. Each hospital's model moves towards its own kind of data, so the models pull
in different directions, and averaging them gives a poor model.
\item \term{QWK (Quadratic Weighted Kappa)} a score from $-1$ to 1 for how well predicted grades
agree with the true grades. 1 = perfect, 0 = no better than guessing. ``Quadratic'' means the
penalty grows with the \emph{square} of the distance: a mistake of 2 grades costs 4 times a mistake
of 1 grade; a mistake of 4 grades costs 16 times. ``Kappa'' means it is corrected for lucky
guessing. It is the standard score for DR grading (used in Kaggle competitions).
\item \term{Test QWK} QWK measured on the \emph{test} images, which are never used for training or
for choosing models, so it is an honest score.
\item \term{Loss function} the number the model tries to make small during training; it measures
how wrong the predictions are.
\item \term{Cross-entropy (CE)} the usual loss for classification. It only looks at the probability
the model gives to the \emph{true} grade, so every wrong grade is treated the same.
\end{terms}

%% ------------------------------------------------------------------
\slide{5}{Example: three very different hospitals}{0:50}
Let me make this concrete with an example. Imagine three hospitals. A big city hospital sees many
patients, but most of their eyes are healthy. An eye specialist hospital sees many severe cases.
And a village clinic has only a few patients, often with early disease. FedAvg gives each hospital
a vote based on how many images it has. So the big city hospital decides almost everything, and the
global model learns that most eyes are healthy. It starts missing the severe cases from the eye
hospital and the early cases from the village. These bars are only an illustration, but we measured
the same effect: in our extreme-skew test, the two smallest clients got only 5.9 percent of FedAvg's
vote.
\begin{terms}
\item \term{Vote = aggregation weight} how much one hospital's model counts when the server
combines the models. All weights add up to 100\%.
\item \term{Illustration vs.\ measurement} the bars on the slide are a picture of the idea. The
5.9\% is a real number from our logged weights.
\end{terms}

%% ------------------------------------------------------------------
\slide{6}{OrdinalFed has two parts}{0:50}
Our method, OrdinalFed, has two parts. The first part works at each hospital. We add a distance
penalty to cross-entropy. The model predicts an expected grade, and we penalise the squared distance
from the true grade. So predicting Grade 0 for a Grade 4 image costs 16 times more than predicting
Grade 3.

The second part works at the server. The server keeps a small, balanced set of 50 images, 10 from
each grade. After each round, it scores every hospital's model on these images using QWK, and gives
better models more weight. The hospitals do not need to change anything, and no extra data is sent.
\begin{terms}
\item \term{OrdinalLoss} our training loss at each hospital:
$\mathcal{L} = (1-\lambda)\,\text{CE} + \lambda\,(\hat{y}-y)^2$. It mixes cross-entropy with a
``distance'' penalty.
\item \term{$\lambda$ (lambda) = 0.5} how much of each part we use. 0.5 means half CE, half
distance penalty.
\item \term{Softmax probabilities} the model gives a probability for each grade, adding up to 100\%.
Example: 10\% Grade 0, 20\% Grade 1, 70\% Grade 2.
\item \term{Expected grade $\hat{y}$} the average grade, weighted by those probabilities:
$\hat{y} = \sum_c c \cdot p_c$. In the example above: $0(0.1)+1(0.2)+2(0.7) = 1.6$.
\item \term{Squared distance $(\hat{y}-y)^2$} if the true grade is 4: predicting 0 costs
$(0-4)^2 = 16$; predicting 3 costs $(3-4)^2 = 1$. So the far mistake costs 16 times more.
\item \term{Class weights $\mathbf{c} = [1, 5, 1.5, 10, 4]$} rare grades count more in CE (Grade 3
is the rarest, so it gets weight 10), so the model does not ignore them.
\item \term{Server validation set $\mathcal{V}$} the 50 images (10 per grade) the server keeps, the
``exam''. Balanced, so every grade counts equally. Taken from the public DDR test split, not from
the hospitals.
\end{terms}

%% ------------------------------------------------------------------
\slide{7}{How the server sets the weights}{0:40}
Here is how the server sets the weights. First, it scores each model on its 50 images. Second, a
softmax turns these scores into quality weights. Third, it mixes two things: each hospital's share
of the images, and its quality weight. A factor called beta controls the mix. In round 1, beta is
0.5, so image counts and quality count equally. Beta then falls smoothly to zero, so in the last
round only quality is used. We never drop a hospital, because a score from only 50 images is noisy.
\begin{terms}
\item \term{$\kappa_k^t$ (kappa)} the QWK of hospital $k$'s model in round $t$, measured on the 50
server images.
\item \term{Floor $\kappa_{\min} = 0.1$} any score below 0.1 is raised to 0.1, so one bad round
cannot give a hospital almost zero weight.
\item \term{Softmax with temperature $\tau = 3$} a formula that turns scores into weights that add up
to 1: $q_k = e^{\kappa_k/\tau} / \sum_j e^{\kappa_j/\tau}$. The \emph{temperature} controls how
``sharp'' it is: a high $\tau$ gives almost equal weights; a low $\tau$ gives the best model most of
the weight. With $\tau = 3$, the biggest weight is at most about 1.35 times the smallest.
\item \term{$\beta(t)$ (beta) and the cosine schedule} the mixing factor:
weight $= \beta \cdot \text{image share} + (1-\beta) \cdot \text{quality}$.
$\beta(t) = 0.5\cos\!\big(\tfrac{\pi(t-1)}{2(T-1)}\big)$: it starts at 0.5 (round 1) and falls
smoothly, like a quarter of a cosine wave, to 0 (round 10). Why: early scores are unreliable, so we
trust image counts first and quality scores later.
\item \term{Minimum weight $\omega_{\min} = 0.03$} a safety rule: every hospital gets at least 3\%.
In our setting it never triggers, because every hospital already gets about 8\% or more.
\item \term{Noisy} a score from only 50 images changes a lot by chance, so it is not fully reliable.
\end{terms}

%% ------------------------------------------------------------------
\slide{8}{Same exam for every hospital}{0:45}
So, back to our three hospitals. OrdinalFed fixes the problem in a simple way. Think of it as giving
every hospital the same exam. The server keeps 50 images, 10 from each grade, and tests every
hospital's model on them. Then the votes slowly move from ``how many images do you have'' to ``how
well did you do on the exam''. And every hospital keeps a minimum vote, so even a tiny village
clinic is never ignored. In our extreme-skew test, the vote of the two smallest hospitals went up
from 5.9 percent to 40.8 percent, and the worst-round QWK went up from 0.378 to 0.514. The hospitals
did not have to change anything.
\begin{terms}
\item \term{Worst-round QWK} the lowest test QWK over the 10 rounds. It shows how badly a method can
fail. Higher worst round = more stable.
\item \term{Logged weights} during training we saved every hospital's weight in every round. The
5.9\% $\rightarrow$ 40.8\% comes from these records (round 10, extreme skew).
\end{terms}

%% ------------------------------------------------------------------
\slide{9}{Experimental setup}{0:35}
We tested on the DDR dataset, which has 12,522 gradable fundus images, split into training,
validation and test sets. We used five hospital clients and an EfficientNet-B3 model at 224 by 224
pixels. Each method ran for 10 rounds with 5 local epochs, and every experiment was repeated with
three random seeds. We compared OrdinalFed with FedAvg, FedProx, an accuracy-weighted baseline
called AccWt, and OrdinalFed combined with FedProx. We also tested every trained model, without
retraining, on a second dataset, APTOS 2019.
\begin{terms}
\item \term{DDR dataset} a public DR dataset collected in China (13,673 photos; 12,522 are good
enough to grade). Split: 6,260 training, 2,503 validation, 3,759 test images.
\item \term{Gradable} the photo is clear enough for a doctor to give a grade.
\item \term{Training / validation / test split} training = to learn; validation = to choose the
best round; test = final, honest score (never used for learning or choosing).
\item \term{EfficientNet-B3} a well-known image network (CNN) that is accurate but not too big
(about 12 million weights). ``B3'' is a medium size in the EfficientNet family.
\item \term{CNN (convolutional neural network)} a neural network made for images; it learns small
patterns (edges, spots) and then bigger ones.
\item \term{224 $\times$ 224 pixels} the size every image is shrunk to. Small images train fast but
can hide tiny lesions.
\item \term{Warm start} all FL runs start from a centralised model trained for 30 epochs (best
validation QWK 0.834 at epoch 23), like a hospital adapting a pre-trained model.
\item \term{Local epoch} one pass through all of a hospital's own images. Each round has 5.
\item \term{Seed} a number that fixes the randomness (shuffling, augmentation). Three seeds = we ran
everything three times with different randomness and report \emph{mean $\pm$ standard deviation}
(average and spread).
\item \term{Training settings} batch size 64 (images per update), Adam optimiser (the update rule),
learning rate $10^{-4}$ (step size). \emph{Augmentation}: random flips, small rotations and colour
changes, so the model does not memorise the images (\emph{overfitting}).
\item \term{FedProx} FedAvg plus a ``rubber band'' at each hospital: the \emph{proximal term}
$\tfrac{\mu}{2}\lVert w - w^t\rVert^2$ pulls each local model back towards the global model, so it
does not drift too far. We used $\mu = 0.01$.
\item \term{AccWt} the same server weighting as ours, but it scores hospitals by plain
\emph{accuracy} instead of QWK, and hospitals use normal cross-entropy. It tests whether QWK itself
matters.
\item \term{OrdinalFed + Prox} our method plus FedProx's rubber band at the hospitals.
\item \term{APTOS 2019 (external validation)} a different public dataset (3,662 photos from India,
from a Kaggle competition). Testing on it, without retraining, checks whether the models work on
new hospitals and cameras.
\end{terms}

%% ------------------------------------------------------------------
\slide{10}{Three ways to split the data}{0:40}
We split the data across the five clients in three ways. Each bar is one client, and the colours
show its mix of grades. In the IID split, on the left, every client gets a random, similar mix.
\act{Point to the middle.} The middle panel is the hard case, a Dirichlet split with alpha 0.1. Here
each client sees mostly one grade. Client 0 has only 32 images, all Grade 0, and Client 4 has 2,087
images, almost all Grade 2. \act{Point right.} The right panel models five real types of clinic, from
screening clinics to specialist hospitals.
\begin{terms}
\item \term{Partition / split} how the training images are divided among the five hospitals.
\item \term{IID split} random and equal: every hospital gets 1,252 images with the same grade mix.
The easy case.
\item \term{Dirichlet split, $\alpha$ (alpha) $= 0.1$} a standard way to make uneven splits for FL
research. $\alpha$ controls how uneven: small $\alpha$ (like 0.1) = each hospital gets mostly one
grade; large $\alpha$ = close to IID. Sizes vary too: 32 to 2,087 images.
\item \term{Real-site split} we designed five clinic types by hand (general screening, rural
primary care, specialist referral, DR clinic, urban hospital) to look like real life.
\end{terms}

%% ------------------------------------------------------------------
\slide{11}{Test QWK after every round}{0:40}
This figure shows the test QWK after each round. Orange is FedAvg, green is FedProx, blue is
OrdinalFed, and the dashed line is the centralised model trained on all the data. On the IID and
real-site splits, on the left and right, all methods stay close to the centralised model.
\act{Point to the middle panel.} The middle panel is the extreme skew. FedAvg collapses from round 5
onwards. OrdinalFed avoids the collapse, but it goes up and down. FedProx stays flat.
\begin{terms}
\item \term{Centralised model} one model trained with all the images in one place (no FL). It is
the reference: test QWK 0.682.
\item \term{Shaded band} $\pm$ one standard deviation over the three seeds: a wide band means the
result changes a lot from run to run.
\end{terms}

%% ------------------------------------------------------------------
\slide{12}{Worst round under extreme skew}{0:35}
Here are the numbers for the worst round under extreme skew, averaged over three seeds. FedAvg drops
to 0.378. OrdinalFed raises the worst round to 0.514, and the mean over all rounds from 0.535 to
0.604. But FedProx, which changes local training, is the most stable method, at 0.641. Adding
FedProx's term to OrdinalFed brings it close, at 0.627, but it does not beat FedProx.
\begin{terms}
\item \term{Stability} how steady a method's score stays across rounds. We measure it with the
\emph{mean}, \emph{worst} and \emph{final} test QWK over the 10 rounds.
\end{terms}

%% ------------------------------------------------------------------
\slide{13}{Why does it work?}{0:50}
Why does OrdinalFed help? We removed one part at a time. The OrdinalLoss alone does not stop the
collapse: 0.380, the same as FedAvg. QWK weighting raises the worst round to 0.586. But simply giving
every hospital the same weight does even better, at 0.630. The logged weights explain this. The two
smallest clients hold 5.9 percent of the images, so FedAvg gives them 5.9 percent of the weight.
OrdinalFed raises their weight to 40.8 percent by round 10. So the real gain comes from giving small
hospitals a bigger voice. With only 50 server images, the QWK scores themselves are too noisy to rank
the hospitals.
\begin{terms}
\item \term{Ablation study} test the method with one part removed or swapped at a time, to see
which part really causes the improvement (like removing one ingredient from a recipe).
\item \term{Image-count weights} FedAvg's weighting (share of images).
\item \term{Equal (uniform) weights $\omega = 1/K$} every hospital gets the same weight: with
$K = 5$ hospitals, 20\% each.
\item \term{Sensitivity analysis} (in the paper, Table III) we also changed one setting at a time
($\lambda = 0.1$, $\beta_0 = 0$, $\tau = 0.5$). All three did a bit better than the default, but
still below equal weights. The settings were fixed before the runs and compared only on validation
data, to stay fair.
\end{terms}

%% ------------------------------------------------------------------
\slide{14}{OrdinalFed or FedProx?}{1:00}
A natural question is: why use OrdinalFed instead of FedProx? The difference is where the fix
happens. FedProx changes the training at every hospital: each one must add an extra term to its
training code and tune a setting. OrdinalFed changes only the server, and the hospitals keep their
normal training. To be honest, FedProx was more stable in our tests: its worst round was 0.641,
against 0.514 for OrdinalFed. But OrdinalFed is useful in three situations. First, when you can
change only the server, for example when small clinics run simple, fixed software. Second, when
hospital sizes are very different and small clinics would be drowned out. Third, when catching rare
severe cases matters: OrdinalFed had the best Grade 3 recall, 0.36 against 0.31. And you do not have
to choose: OrdinalFed plus FedProx's term reached 0.627, so the two work together.
\begin{terms}
\item \term{Client-side vs.\ server-side} client-side methods (FedProx) change how each hospital
trains; server-side methods (OrdinalFed's weighting) change only how the server combines.
\item \term{Tune a setting ($\mu$)} FedProx's strength $\mu$ must be chosen; too small does
nothing, too big stops hospitals from learning.
\item \term{Recall} of all the images that truly are Grade 3, the share the model found. Recall 0.36
= it found 36 out of every 100 Grade 3 images. (\emph{Precision}: of the images the model called
Grade 3, the share that really were. \emph{F1}: one number that balances both.)
\end{terms}

%% ------------------------------------------------------------------
\slide{15}{More results}{0:40}
A few more results. On the IID and realistic multi-site splits, all methods were within 0.01 QWK of
each other, so they perform about the same when the skew is mild. OrdinalFed had the best recall on
Grade 3, the rarest severe grade: 0.36, against 0.27 to 0.31 for the others. On the external APTOS
dataset, all models scored between about 0.60 and 0.66, close to the centralised model's 0.641, so
they transfer well. And the server cost is small: 1.2 seconds per round, with no extra
communication.
\begin{terms}
\item \term{Transfer / generalise} a model trained on one dataset still works on another dataset
from different hospitals and cameras.
\item \term{Overhead} the extra cost of our method: 1.2 s of server time per round (about 0.6\% of a
195-second round on an Intel Arc B580 GPU). \emph{No extra communication}: hospitals send the same
model as in FedAvg.
\end{terms}

%% ------------------------------------------------------------------
\slide{16}{Takeaways}{0:45}
To conclude, there are three takeaways. First, OrdinalFed's server-side weighting prevents most of
FedAvg's collapse under extreme skew, with no change at the hospitals and no extra communication.
Second, the gain comes from giving small hospitals more weight; equal weights work even better.
Third, FedProx is still the most stable method. Our study also has limits: the 50-image server set is
small, no method detects Grade 1 well, and we trained on one dataset at low resolution. In future
work, we will test larger server sets, higher resolution, more datasets, and formal privacy
guarantees.
\begin{terms}
\item \term{Limitation} an honest weak point of the study that future work should fix.
\item \term{Resolution} image size in pixels; higher resolution shows tiny lesions better.
\item \term{Formal privacy guarantees (differential privacy)} a mathematical promise that the shared
model does not reveal any single patient's data. FL alone keeps images at home, but does not give
this promise.
\end{terms}

\slide{17}{Thank you}{0:05}
Thank you for listening. I am happy to take your questions.

%% ------------------------------------------------------------------
\newpage
{\color{navy}\Large\bfseries Quick glossary (A--Z)}\par\smallskip
{\small
\begin{longtable}{@{}>{\bfseries}p{3.6cm}p{11.1cm}>{\leavevmode\color{muted}}p{0.9cm}@{}}
\toprule
Term & Easy meaning & Slide\\ \midrule\endhead
Ablation & Remove one part at a time to see what really helps. & 13\\
AccWt & Our weighting, but scored by accuracy instead of QWK. & 9\\
Aggregation & Combining hospital models into one. & 1, 3\\
APTOS 2019 & Second dataset (India), used only for testing. & 9\\
$\beta$ (beta) & Mix of image share and quality; 0.5 $\rightarrow$ 0 over the rounds. & 7\\
Centralised model & One model trained on all images in one place; the reference. & 11\\
Client & One hospital in FL. & 3\\
Client drift & Local models move towards their own data and apart from each other. & 4\\
Collapse & The global model suddenly gets much worse. & 4\\
Cross-entropy (CE) & Usual loss; treats every wrong grade the same. & 4\\
DDR & Main dataset (China), 12,522 gradable images. & 9\\
Dirichlet split ($\alpha$) & Uneven split; small $\alpha$ = each hospital mostly one grade. & 10\\
EfficientNet-B3 & Medium-size image network (CNN). & 9\\
Epoch (local) & One pass through a hospital's own images. & 9\\
Expected grade $\hat{y}$ & Probability-weighted average grade. & 6\\
FedAvg & Standard FL; weights hospitals by number of images. & 3\\
FedProx & FedAvg + ``rubber band'' (proximal term) at each hospital. & 9, 14\\
IID / non-IID & Hospitals' data alike / different. & 4, 10\\
$\kappa$ (kappa) & A hospital's QWK on the 50 server images. & 7\\
$\lambda$ (lambda) & Balance between CE and the distance penalty (0.5). & 6\\
Label skew & Hospitals have different mixes of grades. & 4\\
OrdinalLoss & CE + squared-distance penalty. & 6\\
Proximal term ($\mu$) & FedProx's pull towards the global model; $\mu = 0.01$. & 9\\
QWK & Agreement score for ordered grades; 1 = perfect. & 4\\
Recall & Share of true cases the model finds. & 14\\
Round & One send--train--return--combine cycle. & 3\\
Seed & Number fixing randomness; we used 3. & 9\\
Server set $\mathcal{V}$ & The 50 ``exam'' images, 10 per grade. & 6\\
Softmax & Turns numbers into probabilities/weights that add up to 1. & 6, 7\\
Temperature $\tau$ & How sharp the softmax is; $\tau = 3$ = gentle. & 7\\
Test split & Images used only for the final, honest score. & 9\\
Validation split & Images used to choose the best round. & 9\\
\bottomrule
\end{longtable}}

%% ------------------------------------------------------------------
{\color{navy}\Large\bfseries Harder questions and easy answers}
\begin{enumerate}[leftmargin=*,itemsep=6pt]
  \item \textbf{Why QWK and not accuracy?}\\
  Accuracy only counts right or wrong. QWK also counts \emph{how} wrong: calling Grade 4 ``Grade 0''
  is punished much more than ``Grade 3''. For ordered grades, that matters, and QWK is the standard DR
  score.
  \item \textbf{Why does QWK weighting not beat equal weighting?}\\
  The scores come from only 50 images, so they are noisy, and with $\tau = 3$ the weights end up
  almost equal anyway. The real benefit is that small hospitals get more weight, which equal weights
  give directly.
  \item \textbf{Why $\tau = 3$ and $\beta_0 = 0.5$?}\\
  They are the settings our method was designed with. We also tested other values, fixed before the
  runs: a sharper $\tau = 0.5$ and $\beta_0 = 0$ did a bit better on validation (0.795 and 0.793 vs.\
  0.761), but still below equal weights (0.810).
  \item \textbf{Why a cosine schedule for $\beta$?}\\
  Early in training the scores are unreliable, so we trust image counts first and move smoothly to
  quality scores. A cosine curve gives a smooth, gentle change.
  \item \textbf{Table I shows OrdinalFed lowest under $\alpha = 0.1$ (0.640). Why?}\\
  Table I uses the round with the best \emph{validation} score. For FedAvg that is always round 1,
  before it collapses, so Table I hides the collapse. That is why we also report mean, worst and
  final rounds (Table II).
  \item \textbf{Is keeping 50 labelled images on the server a privacy problem?}\\
  In our tests they come from the public DDR test split, not from the hospitals. In real use, the
  server needs a small labelled set; this is a limitation.
  \item \textbf{Why is test QWK (0.682) lower than validation QWK (0.834)?}\\
  The test split has a different mix of cameras, and many test Grade 2 images are predicted as
  Grade 0 (about 45\%, against 12\% in validation). This affects all methods the same way, so we
  report test results.
  \item \textbf{Why is Grade 1 so hard?}\\
  Grade 1 is only 5\% of the training images, and its signs (tiny microaneurysms) are hard to see at
  224 $\times$ 224 pixels. Higher resolution is future work.
  \item \textbf{What does the OrdinalLoss add, then?}\\
  Ordinal losses help in centralised training. In our federated tests it gave no measurable benefit
  (and slightly lowered results). We report that honestly.
  \item \textbf{Did you tune FedProx's $\mu$?}\\
  No. We used $\mu = 0.01$, a common value, the same for all FedProx runs.
  \item \textbf{Why only five hospitals and ten rounds?}\\
  To keep all experiments (three seeds, three splits, ablations) affordable on one GPU. More clients
  and longer training are future work.
\end{enumerate}

\end{document}
